医学图像分割作为医学影像分析中的核心任务，在疾病诊断、手术规划、病灶评估及疗效跟踪等临床环节中具有不可替代的作用。然而，当前主流基于深度学习的医学图像分割方法通常依赖于大量高质量、精细标注的训练数据。在实际临床环境中，获取这样的标注数据面临诸多挑战：标注过程耗时耗力、需依赖专业医师知识、不同医疗机构间数据分布差异大等问题，导致标注成本高昂且数据规模受限。此外，现有模型往往在特定数据集上表现良好，而在跨机构、跨模态、跨解剖部位的泛化能力较差，严重制约了其在实际医疗场景中的推广应用。

因此，本研究旨在针对标注稀缺这一现实瓶颈，系统构建一套具有强泛化能力的自适应医学图像分割理论与方法体系。具体而言，本研究首先基于CA数据集，构建了认知不确定性引导的半监督学习框架，旨在Mean Teacher架构中通过量化模型预测的不确定性来提升伪标签质量，从而改善模型性能。随后，在网络结构设计方面，深入探究了U-Net与Transformer的不同融合范式，并进一步在跳跃连接中引入通道注意力机制，以增强模型对多尺度特征的自适应筛选与整合能力。